\chapter{Input Formats and Conventions}
\label{ch:input}

OpenSG-2.0 reads three kinds of structure-gene (SG) description, plus one
laminate shorthand and one legacy converter. Which dialect you write is decided
by the geometry you have, not by the answer you want: a thin-walled cross-section
contour is a 1-D shell SG, a sheet meshed as a surface in space is a 3-D shell
SG, and anything meshed as a filled domain (a laminate through its thickness, a
honeycomb unit cell, a lattice block) is a solid SG. The five dialects are listed
in Table~\ref{tab:input-dialects}.

\begin{table}[htbp]
\centering
\small
\caption{The five input dialects, the package that reads each, and the example
file shipped with the code.}
\label{tab:input-dialects}
\begin{tabular}{@{}lll@{}}
\toprule
dialect & read by & example file \\
\midrule
1-D shell SG (cross-section contour) & \texttt{opensg\_shell}
  & \texttt{iea\_s10\_shell.yaml} \\
3-D shell SG (surface in space) & \texttt{opensg\_shell}
  & \texttt{schwarz\_p\_3Dshell.yaml} \\
solid SG (1-D / 2-D / 3-D mesh) & \texttt{opensg\_solid}
  & \texttt{RHC\_SW\_2UC\_45.yaml} \\
SwiftComp \texttt{.sc} & converted to the solid SG YAML
  & \texttt{RHC\_SW\_2UC\_45.sc} \\
through-thickness layup database & \texttt{opensg\_solid.rm\_plate\_1D}
  & \texttt{layup\_db.yaml} \\
\bottomrule
\end{tabular}
\end{table}

The two shell dialects share one schema --- the same six top-level keys,
differing only in element dimension --- while the solid dialect is a separate,
deliberately minimal schema whose material blocks mirror SwiftComp's.

Every run, in every dialect, writes a timed \texttt{.out} in the SwiftComp
\texttt{.K} layout. The writer is \texttt{opensg\_solid.sg\_homo.write\_sc\_K}
and it always emits, in order: an \texttt{ OpenSG} banner line naming the model,
\texttt{The Effective <Name> Stiffness Matrix}, \texttt{The Effective <Name>
Compliance Matrix}, and a \texttt{ Time taken: <s> sec} footer. The
\texttt{<Name>} infix is \texttt{Timoshenko} for a beam $6\times6$,
\texttt{Classical Plate} for a $6\times6$ ABD, \texttt{Reissner-Mindlin Plate}
for an $8\times8$ ABDG, and empty for the equivalent 3-D solid law --- which
alone is written with \texttt{constants=True} and therefore also prints
\texttt{The Engineering Constants (Approximated as Orthotropic)}. The full
contract is \texttt{Rules/output\_format\_and\_timing.md} at the repository
root.

%========================================================================
\section{The 1-D shell SG YAML}
\label{sec:input-shell1d}

This is the cross-section of a thin-walled beam: the wall midline (or mold line)
meshed with 2-node line segments, each segment carrying a laminate. It is the
input to \texttt{opensg\_shell.build\_rm\_bundle}, which returns the Timoshenko
$6\times6$.

The example \texttt{iea\_s10\_shell.yaml} in
\texttt{examples/OpenSG\_shell/1\_get\_beam\_props\_from\_shell\_cross\_section/}
is station $r/R = 0.2$ of the IEA-22 blade: 307 nodes, 310 line elements, 6
layups, 6 materials. Its top-level keys, in the order OpenSG's own writers emit
them, are listed in Table~\ref{tab:input-shell-keys}.

\begin{table}[htbp]
\centering
\small
\caption{Top-level keys of the 1-D shell SG YAML.}
\label{tab:input-shell-keys}
\begin{tabular}{@{}lp{10.2cm}@{}}
\toprule
key & content \\
\midrule
\texttt{nodes} & one bracketed, space-separated triple per node:
  the two cross-section coordinates then the beam-axis coordinate \\
\texttt{elements} & one bracketed, space-separated pair per line element,
  1-based node ids \\
\texttt{sets} & \texttt{element:} a list of \texttt{\{name, labels\}} groups;
  \texttt{labels} are 1-based element indices \\
\texttt{sections} & one entry per layup, tied to a set through
  \texttt{elementSet}, carrying the \texttt{layup} ply table \\
\texttt{elementOrientations} & one row of nine direction cosines per element,
  a genuine YAML list of floats \\
\texttt{materials} & named orthotropic materials with \texttt{density} and an
  \texttt{elastic} block of \texttt{E}, \texttt{G}, \texttt{nu} triples \\
\texttt{reference} & one scalar naming the surface the wall laws are referenced
  to: \texttt{center}, \texttt{oml}, \texttt{oml\_flip} or \texttt{iml} \\
\bottomrule
\end{tabular}
\end{table}

\subsection{\texttt{nodes}}

One row per node, three coordinates. In the shell dialect the row is written as a
bracketed, \emph{space-separated} triple --- a single YAML scalar inside a flow
sequence --- and not as a comma list:

\begin{lstlisting}[style=shellst,caption={The first two node rows of iea\_s10\_shell.yaml.}]
nodes:
- [4.70852238 0.09722431 0.00000000]
- [4.64563825 0.11238283 0.00000000]
\end{lstlisting}

The space-separated form is required, not merely conventional: the orientation
plotter \texttt{fe\_jax/orient\_plot.py} parses each row with
\texttt{str(row[0]).split()}, so a comma-separated list of floats will not plot.
The loader \texttt{sg\_mesh.load\_ring\_ref} is more forgiving --- its
\texttt{\_row} helper replaces commas by spaces before splitting --- but writing
the space form keeps every consumer happy at once.

The component order is fixed and is the source of most first-run confusion. The
loader sets \texttt{ax = 2} and \texttt{cross = [0, 1]}, meaning the
\textbf{third column is the beam axis} and the first two columns are the
in-plane cross-section pair $(y_2, y_3)$. For a plane cross-section the beam-axis
column is identically zero, exactly as in the excerpt above. If you are writing a
cross-section by hand and your third column is not all zeros, you have written a
segment, not a cross-section.

\subsection{\texttt{elements}}

One row per element, the node ids of a 2-node line segment, 1-based, again
space-separated inside brackets:

\begin{lstlisting}[style=shellst,caption={The first two element rows; element k here is element label k everywhere else in the file.}]
elements:
- [1 2]
- [2 3]
\end{lstlisting}

The loader detects the base with \texttt{if cells.min() == 1: cells = cells - 1},
so a 0-based mesh is also accepted, but every generator in the repository writes
1-based. Element $k$ in this list is element label $k$ in \texttt{sets}, and
row $k$ of \texttt{elementOrientations} belongs to it. The contour need not be a
single loop: the last element of the IEA station is \texttt{- [307 65]}, which
closes a shear web back onto a skin node.

\subsection{\texttt{sets} and \texttt{sections}}

\texttt{sets} groups element labels by name; \texttt{sections} gives each named
group its laminate. The tie between them is the \texttt{elementSet} field, which
must equal a set \texttt{name}:

\begin{lstlisting}[style=shellst,caption={Abridged sets and sections blocks. In the real file layup\_0 lists 25 labels (elements 1-12 and 209-221) and there are six sets and six sections.}]
sets:
  element:
  - name: layup_0
    labels:
    - 1
    - 2
    - 3
sections:
- type: shell
  elementSet: layup_0
  layup:
  - - gelcoat
    - 0.0005
    - 0.0
  - - glass_triax
    - 0.00300294
    - 0.0
  - - glass_uniax
    - 0.02609467
    - 0.0
  - - glass_triax
    - 0.00300294
    - 0.0
\end{lstlisting}

Labels are 1-based element indices into \texttt{elements}.
\texttt{sg\_materials.\_material\_by\_section} builds one ABD per entry of
\texttt{sections} and keys it by the section's \emph{position in the list}; the
loader then maps every element through \texttt{rsub[label - 1] =
section\_index}. An element that appears in no set silently falls to section 0,
so every element should appear in exactly one set. \texttt{type: shell} is
descriptive metadata --- the shell readers key only off \texttt{elementSet} and
\texttt{layup}.

Each \texttt{layup} row is \texttt{[material, thickness, angle\_deg]}: a
material name that must exist in \texttt{materials}, a ply thickness in metres,
and a ply angle in degrees. The list is ordered \textbf{outer face first} and
stacks inward along $e_3$: \texttt{compute\_ABD\_matrix} documents its input as
layer thicknesses bottom to top with the reference surface at that bottom
(outer) face, and \texttt{emit\_station\_abd} records the same order in its
\texttt{convention} string as \texttt{plies OML->IML}. The sum of the row
thicknesses is the wall thickness used for the reference shift; for
\texttt{layup\_0} above that is $0.0005 + 0.00300294 + 0.02609467 + 0.00300294 =
0.03260055$~m, which is the \texttt{thickness} the cached ABD file reports for
that layup.

\subsection{\texttt{elementOrientations}}

One row of nine direction cosines per element, $[e_1(3)\;e_2(3)\;e_3(3)]$, in the
same component order as the node rows. Unlike \texttt{nodes} and
\texttt{elements}, these rows are true YAML lists of floats:

\begin{lstlisting}[style=shellst,caption={One orientation row: e1 = [0,0,1] out of plane, then the in-plane e2 and e3.}]
elementOrientations:
- [0.0, 0.0, 1.0, -0.9721541292349604, 0.2343423756204075, 0.0, -0.2343423756204075,
  -0.9721541292349604, 0.0]
\end{lstlisting}

The meaning of the triad is covered in Section~\ref{sec:input-orientation}.

\subsection{\texttt{materials}}

A list of named materials, each with a \texttt{density} and an \texttt{elastic}
block holding three-component \texttt{E}, \texttt{G} and \texttt{nu}:

\begin{lstlisting}[style=shellst,caption={One material block from iea\_s10\_shell.yaml.}]
materials:
- name: glass_triax
  density: 1940.0
  elastic:
    E:
    - 28211400000.0
    - 16238800000.0
    - 15835500000.0
    G:
    - 8248220000.0
    - 3491240000.0
    - 3491240000.0
    nu:
    - 0.497511
    - 0.18091
    - 0.27481
\end{lstlisting}

The triples are $[E_1, E_2, E_3]$, $[G_{12}, G_{13}, G_{23}]$ and $[\nu_{12},
\nu_{13}, \nu_{23}]$ --- the argument order of
\texttt{fe\_jax.msg\_materials.build\_stiffness\_6x6}, which fills the
orthotropic compliance
\[
S_{11} = 1/E_1,\quad S_{12} = -\nu_{12}/E_1,\quad S_{13} = -\nu_{13}/E_1,\quad
S_{23} = -\nu_{23}/E_2,
\]
\[
S_{44} = 1/G_{23},\quad S_{55} = 1/G_{13},\quad S_{66} = 1/G_{12},
\]
and inverts it to the $6\times6$ stiffness. Units follow the input; the examples
are SI, so moduli are in Pa, thicknesses in m and \texttt{density} in
kg/m$^3$. \texttt{density} feeds the section mass per unit area that
\texttt{emit\_station\_abd} reports as \texttt{mass\_per\_area} (62.892567
kg/m$^2$ for \texttt{layup\_0} of this station); it is not needed for the
stiffness.

\subsection{\texttt{reference}}

\begin{lstlisting}[style=shellst,caption={The whole reference block -- one scalar, the last line of the file.}]
reference: center
\end{lstlisting}

This single scalar names the surface the wall ABD --- and therefore the whole
cross-section model --- is referenced to, and it is the single source of truth.
\texttt{build\_rm\_bundle(shell\_yaml)} reads it when its \texttt{ref} argument
is left \texttt{None} via \texttt{ref = d.get("reference", "center")}, and the
value is immediately turned into a thickness fraction
\texttt{frac = \{"center": 0.5, "oml": 0.0, "oml\_flip": 1.0, "iml": 1.0\}} that
then drives the ring laminate reference, the plate-SG \texttt{z\_ref}, the
emitted ABD file, and the depth conversion used in stress recovery. Absent, it
defaults to \texttt{center}. The four accepted values and the fraction each maps
to are listed in Table~\ref{tab:input-reference}.

\begin{table}[htbp]
\centering
\small
\caption{Values of \texttt{reference} and the thickness fraction each maps to.}
\label{tab:input-reference}
\begin{tabular}{@{}llc@{}}
\toprule
value & meaning & fraction from the outer face \\
\midrule
\texttt{center} & contour is the laminate mid-surface; ABD shifted by $t/2$ & 0.5 \\
\texttt{oml} & contour is the outer mold line; laminate stacks inward, no shift & 0.0 \\
\texttt{oml\_flip} & diagnostic: reference moved the full thickness the other way & 1.0 \\
\texttt{iml} & contour is the inner mold line; laminate stacks outward & 1.0 \\
\bottomrule
\end{tabular}
\end{table}

Mid-surface meshes (tubes, ellipses, generated cross-sections) use
\texttt{center}; airfoil YAMLs written on the OML use \texttt{oml}. Getting this
wrong shifts $\mathbf{B}$ and $\mathbf{D}$ by a parallel-axis term of the wall
thickness and moves every bending and torsion entry of the beam $6\times6$, which
is why it is recorded in the file rather than passed at the call site.

\subsection{Running it}

\begin{lstlisting}[style=shellst]
cd examples/OpenSG_shell/1_get_beam_props_from_shell_cross_section
python beam_homo_shell.py
\end{lstlisting}

The driver sets \texttt{name = "iea\_s10"}, calls
\texttt{build\_rm\_bundle(name + "\_shell.yaml")}, and prints the $6\times6$
with the diagonal labelled \texttt{EA GA2 GA3 GJ EI2 EI3}:

\begin{lstlisting}[style=opensg,caption={The whole compute part of beam\_homo\_shell.py.}]
from opensg_shell import build_rm_bundle, auto_emit

name = "iea_s10"
B = build_rm_bundle(name + "_shell.yaml")
C6 = np.asarray(B["Timo"])
LBL = ["EA", "GA2", "GA3", "GJ", "EI2", "EI3"]
print("diagonal: " + "  ".join("%s=%.5g" % (LBL[i], C6[i, i]) for i in range(6)))

auto_emit(name + "_shell.yaml", out_png=name + "_orient.png")
\end{lstlisting}

Four files come back. \texttt{iea\_s10\_shell\_Timo.out} holds the Timoshenko
$6\times6$ and its compliance under the banner \texttt{ OpenSG msg-shell beam
model [ext sh2 sh3 twist bend2 bend3]}; for this station the diagonal reads
$EA = 2.766\times10^{10}$~N, $GA_2 = 7.186\times10^{8}$~N,
$GA_3 = 4.217\times10^{8}$~N, $GJ = 2.438\times10^{9}$~N\,m$^2$,
$EI_2 = 3.514\times10^{10}$~N\,m$^2$, $EI_3 = 6.813\times10^{10}$~N\,m$^2$, and
the footer records \texttt{ Time taken: 8.02 sec}.
\texttt{iea\_s10\_shell\_ABDG.out} holds one $8\times8$ RM wall law per section,
written by \texttt{sg\_homo.write\_abdg\_out} with the header
\texttt{rows [eps11 eps22 2eps12 | K11 K22 K12+K21 | 2g13 2g23]}, so you can read
off, for instance, the $A_{11} = 1.370\times10^{9}$~N/m and the transverse-shear
$G_{11} = 8.283\times10^{7}$~N/m of \texttt{layup\_0} directly.
\texttt{iea\_s10\_mesh.png} is the contour coloured by layup and
\texttt{iea\_s10\_orient.png} is the orientation check discussed in
Section~\ref{sec:input-orientation}. On the first run \texttt{build\_rm\_bundle}
additionally caches the per-station wall laws as
\texttt{abd/iea\_s10\_shell\_abd.yaml} at the same reference the YAML declared
(here written out as \texttt{reference: mid}); the write is guarded by an
existence test on that path, so later runs reuse the cached file rather than
recompute it.

%========================================================================
\section{The 3-D shell SG YAML}
\label{sec:input-shell3d}

Same six keys, but the mesh is a surface embedded in 3-D: \texttt{nodes} are
genuine three-dimensional points and \texttt{elements} are 3-noded triangles or
4-noded quads. There is no \texttt{reference} key --- \texttt{shell\_sg3d} calls
\texttt{\_material\_by\_section} with \texttt{center\_ref=True} unconditionally,
so every wall law is referenced to its mid-surface. This is the input to
\texttt{opensg\_shell.sg\_homo.shell\_sg3d}, which returns the equivalent 3-D
solid $6\times6$.

\texttt{schwarz\_p\_3Dshell.yaml} is a Schwarz-P triply-periodic minimal-surface
cell: 13536 nodes and 26360 triangles, one aluminium layup at $T = 0.036457$~m
with $E = 69$~GPa and $\nu = 0.3$. Triangles are read and padded to the quad
connectivity by repeating the last node,
\texttt{el = np.array([e + [e[-1]]*(4 - len(e)) for e in el], int) - 1}, so tri
and quad meshes go through one element path; ids are 1-based as in the 1-D
dialect. The mesh itself is shown in Figure~\ref{fig:input-schwarz-mesh}.

\begin{figure}[htbp]
\centering
\includegraphics[width=0.62\textwidth]{\FIGDIR/schwarz_p_mesh.png}
\caption{The Schwarz-P TPMS 3-D shell SG: 13536 nodes, 26360 triangles, one
aluminium layup. The mesh figure is drawn by the driver from the same YAML the
solver reads.}
\label{fig:input-schwarz-mesh}
\end{figure}

\texttt{make\_schwarz\_yaml.py} in the same folder is the canonical way to write
one, and its emission loop shows every row format at once. The frame is built per
triangle from the geometry: $e_3$ is the unit facet normal, $e_2$ is the unit
$p_1 \to p_2$ edge, and $e_1 = e_2 \times e_3$.

\begin{lstlisting}[style=opensg,caption={The emission loop of make\_schwarz\_yaml.py.}]
o = ["nodes:"]
o += ["- [%.12f %.12f %.12f]" % tuple(p) for p in nd]
o.append("elements:")
o += ["- [%d %d %d]" % tuple(e) for e in el]
o.append("elementOrientations:")
for k in range(ne):
    e1 = np.cross(e2[k], e3[k])
    o.append("- [%.6f, %.6f, %.6f, %.6f, %.6f, %.6f, %.6f, %.6f, %.6f]"
             % (*e1, *e2[k], *e3[k]))
o += ["sections:", "- type: shell", "  elementSet: layup_0", "  layup:",
      "  - - alu", "    - %.6e" % T, "    - 0.0",
      "materials:", "- name: alu", "  density: 2700.0", "  elastic:",
      "    E: [%.6e, %.6e, %.6e]" % (E, E, E),
      "    G: [%.6e, %.6e, %.6e]" % (G, G, G),
      "    nu: [%.6f, %.6f, %.6f]" % (nu, nu, nu),
      "sets:", "  element:", "  - name: layup_0", "    labels:"]
o += ["    - %d" % (e+1) for e in range(ne)]
open("schwarz_p_3Dshell.yaml", "w").write("\n".join(o) + "\n")
\end{lstlisting}

Note the asymmetry the loop makes explicit and that trips up hand-written files:
\texttt{nodes} and \texttt{elements} rows are space-separated inside brackets,
\texttt{elementOrientations} rows are comma-separated, and the \texttt{sets}
labels are 1-based. The generator prints the surface area and the relative
density $\text{area} \times T$ so the shell idealisation can be checked against
the solid cell it replaces.

\begin{lstlisting}[style=shellst]
cd examples/OpenSG_shell/4_get_solid_props_from_shell_3D_SG
python make_schwarz_yaml.py
python schwarz_solid_props.py
\end{lstlisting}

The second command writes \texttt{schwarz\_p\_3Dshell\_C3D.out}, whose banner
records the mesh size, the junction-edge count, the boundary treatment and the
bounding-box volume the law is normalized by. For this cell it reads
\texttt{ OpenSG msg-shell equivalent 3D solid (3-D shell SG, 13536 nodes, 26360
elems, 0 junction edges, periodic in 3 dirs, per unit cell 1)}. Because the 3-D
solid law is written with \texttt{constants=True}, the file ends with the
orthotropic-approximated engineering constants --- here
$E_1 = 9.3296\times10^{8}$, $E_2 = 9.3314\times10^{8}$,
$E_3 = 9.3285\times10^{8}$~Pa and $G_{12} = 1.0409\times10^{9}$~Pa --- followed
by \texttt{ Time taken: 66.43 sec}. The near-equality of the three Young's
moduli is the physical check: a Schwarz-P cell is cubic-symmetric, so any
appreciable spread means the mesh or the frames are wrong.

%========================================================================
\section{The solid SG YAML}
\label{sec:input-solid}

The solid dialect describes a meshed \emph{domain} rather than a laminated
\emph{surface}, so it carries no layups and no orientations by default --- the
material blocks are already $6\times6$ stiffnesses or engineering constants.
\texttt{RHC\_SW\_2UC\_45.yaml} is a two-unit-cell honeycomb plate SG with 4251
nodes, 6640 triangles and 3 materials (Figure~\ref{fig:input-rhc-mesh}). The
excerpt below keeps two rows per list and two of the three material blocks; its
keys are collected in Table~\ref{tab:input-solid-keys} and its material types in
Table~\ref{tab:input-mattypes}.

\begin{lstlisting}[style=shellst,caption={Abridged RHC\_SW\_2UC\_45.yaml.}]
dim: 2
scale: 35.248
nodes:
- [17.6240005, 23.3500004, 0.0]
- [-17.6240005, 23.3500004, 0.0]
cells:
- [0, 66, 345]
- [345, 3, 0]
mat_id: [1, 1, 1, 1, 1]
materials:
  1:
    type: 2
    aux: [0.0, 0.0]
    C:
    - [35696.6, 27696.6, 3251.1, 0.0, 0.0, -25368.7]
    - [27696.6, 35696.6, 3251.1, 0.0, 0.0, -25368.7]
    - [3251.1, 3251.1, 8917.8, 0.0, 0.0, -487.1]
    - [0.0, 0.0, 0.0, 3500.0, -500.0, 0.0]
    - [0.0, 0.0, 0.0, -500.0, 3500.0, 0.0]
    - [-25368.7, -25368.7, -487.1, 0.0, 0.0, 27958.5]
  3:
    type: 0
    aux: [0.0, 0.0]
    E: 69000.0
    nu: 0.3
\end{lstlisting}

\begin{table}[htbp]
\centering
\small
\caption{Keys of the solid SG YAML, as consumed by
\texttt{opensg\_solid.sg\_mesh.load\_sg\_input}.}
\label{tab:input-solid-keys}
\begin{tabular}{@{}lp{10.4cm}@{}}
\toprule
key & meaning \\
\midrule
\texttt{dim} & SG spatial dimension, 1, 2 or 3. Only the first \texttt{dim} node
  coordinates are used (\texttt{points = nodes[:, 0:n\_sg]}), so the trailing
  zeros above are padding. \\
\texttt{scale} & the normalization value carried over from the \texttt{.sc}
  trailing line. It is parsed and preserved on the input dict, but the
  homogenizer computes its own SG measure $\omega$ from the mesh and prints that
  in the banner. \\
\texttt{nodes} & comma-separated coordinate rows, always three components. \\
\texttt{cells} & element connectivity, \textbf{0-based} in this dialect (note
  the \texttt{0} in the first row). All elements in one SG must have the same
  node count; a mixed batch raises \texttt{mixed element types in one SG are
  unsupported}. \\
\texttt{mat\_id} & one material id per cell, \textbf{1-based} (the reader forms
  \texttt{mat\_id - 1} internally). \\
\texttt{materials} & a mapping from that id to a material block. \\
\bottomrule
\end{tabular}
\end{table}

Each material block carries a \texttt{type} that selects how the $6\times6$ is
built, matching the SwiftComp material types.

\begin{table}[htbp]
\centering
\small
\caption{Material types in the solid SG YAML and the \texttt{.sc} it mirrors.}
\label{tab:input-mattypes}
\begin{tabular}{@{}clp{7.6cm}@{}}
\toprule
\texttt{type} & fields & interpretation \\
\midrule
0 & \texttt{E}, \texttt{nu} & isotropic; $G = E/[2(1+\nu)]$ \\
1 & \texttt{engineering} & the nine constants
  $E_1\,E_2\,E_3\,G_{12}\,G_{13}\,G_{23}\,\nu_{12}\,\nu_{13}\,\nu_{23}$ \\
2 & \texttt{C} & the $6\times6$ stiffness stored directly --- typically an
  already-rotated ply, so no \texttt{angles} should be supplied with it \\
\bottomrule
\end{tabular}
\end{table}

\texttt{aux} is the auxiliary line of the material block, carried through
verbatim from the \texttt{.sc}; the elastic path does not consume it.

\begin{figure}[htbp]
\centering
\includegraphics[width=0.72\textwidth]{\FIGDIR/rhc_2dsg_mesh.png}
\caption{The RHC honeycomb 2-D solid SG, elements coloured by material id. This
figure is written automatically as \texttt{RHC\_SW\_2UC\_45\_mesh.png} by
\texttt{plot\_sg\_mesh} on every run with \texttt{plot=True} (the default), and
is drawn from the actual parsed mesh.}
\label{fig:input-rhc-mesh}
\end{figure}

The driver overrides the pre-rotated \texttt{type: 2} blocks with engineering
constants plus explicit per-material angles, which is the recommended pattern
when you want the ply angle to appear in exactly one place:

\begin{lstlisting}[style=opensg,caption={plate\_homo\_2dsg.py, the whole user input and the one call.}]
import jax.numpy as jnp
from opensg_solid.sg_homo import plate_homo_2d

n_model = 2                    # 1: Beam; 2: Plate; 3: 3D elastic
name = "RHC_SW_2UC_45"         # reads <name>.yaml

material_param = jnp.array([
    (108e3, 8e3, 8e3, 4e3, 4e3, 3e3, 0.32, 0.32, 0.30),      # ply +45
    (108e3, 8e3, 8e3, 4e3, 4e3, 3e3, 0.32, 0.32, 0.30),      # ply -45
    (69e3, 69e3, 69e3, 26.54e3, 26.54e3, 26.54e3, 0.30, 0.30, 0.30)])
angles = jnp.array([45.0, -45.0, 0.0])   # deg per material; 0.0 = none

r = plate_homo_2d(name + ".yaml", material_param=material_param,
                  angles=angles, n_model=n_model)
print(r["C_eff"])
\end{lstlisting}

\begin{lstlisting}[style=shellst]
cd examples/OpenSG-solid/3_get_plate_props_from_2D_SG
python plate_homo_2dsg.py
\end{lstlisting}

\texttt{n\_model} selects the macro model: 1 beam (Timoshenko), 2 plate, 3
equivalent 3-D solid. The run writes \texttt{RHC\_SW\_2UC\_45.out} ---
\texttt{The Effective Classical Plate Stiffness Matrix}, rows
\texttt{[N11 N22 N12 M11 M22 M12]}, under the banner
\texttt{ OpenSG msg-solid plate model, omega 35.248001, periodic} --- plus
\texttt{RHC\_SW\_2UC\_45\_mesh.png}. The printed $\omega = 35.248001$ is the
in-plane span measured from the mesh; it agrees with the file's
\texttt{scale: 35.248} to six digits, and that agreement is the check that the
mesh and the \texttt{.sc} header describe the same cell. For this run
$A_{11} = 3.3586\times10^{5}$, $D_{11} = 9.0368\times10^{7}$ in the file's own
units, and the footer reads \texttt{ Time taken: 7.83 sec}.

%========================================================================
\section{The SwiftComp \texttt{.sc} input and the converter}
\label{sec:input-sc}

A \texttt{.sc} file can be used directly --- \texttt{load\_sg\_input} dispatches
on the extension and converts --- or converted once and kept as YAML. The
anatomy the reader expects is documented in the
\texttt{opensg\_solid.rm\_plate\_1D.helper.sc\_to\_yaml} module docstring,
implemented in its \texttt{read\_sc}, and summarized in
Table~\ref{tab:input-sc}.

\begin{table}[htbp]
\centering
\small
\caption{Blocks of a SwiftComp \texttt{.sc} file, in file order.}
\label{tab:input-sc}
\begin{tabular}{@{}lp{10.2cm}@{}}
\toprule
block & content \\
\midrule
model header & 1 to 3 short lines (submodel flags and similar) \\
META line & \texttt{dim  n\_nodes  n\_elems  n\_mats} followed by further
  integers; located as the first line holding at least five integers \\
node records & \texttt{n\_nodes} lines of \texttt{id  x [y [z]]} \\
element records & \texttt{n\_elems} lines of \texttt{id  mat\_id  conn}, the
  connectivity zero-padded to a fixed slot count; the padding zeros are not
  nodes \\
material blocks & per material: a header \texttt{mat\_id  mat\_type  n}, an
  auxiliary line such as \texttt{0 0}, then the properties --- type 0 is
  \texttt{E  nu}, type 1 is \texttt{E1 E2 E3 / G12 G13 G23 / v12 v13 v23}, type
  2 is the 21 upper-triangular constants of the $6\times6$ over six lines \\
trailing line & the \texttt{scale} (volume or thickness normalization) \\
\bottomrule
\end{tabular}
\end{table}

In \texttt{RHC\_SW\_2UC\_45.sc} the META line is \texttt{2 4251 6640 3 0 0} --- a
2-D SG with 4251 nodes, 6640 elements and 3 materials --- and the file ends with
the lone line \texttt{35.248}. The file's last triangle, element 6640 of material
3 on nodes 4251/2988/3079, shows the zero padding in full:

\begin{lstlisting}[style=shellst]
6640 3 4251 2988 3079 0 0 0 0 0 0 0 0 0 0 0 0 0 0 0 0 0
\end{lstlisting}

The documented way to run the conversion is the module's own \texttt{convert}:

\begin{lstlisting}[style=opensg,caption={Converting a .sc by hand.}]
from opensg_solid.rm_plate_1D.helper.sc_to_yaml import convert

sc = convert("RHC_SW_2UC_45.sc")
print(sc["dim"], len(sc["nodes"]), len(sc["cells"]), sc["scale"])
\end{lstlisting}

\texttt{convert} writes \textbf{two} files next to the input ---
\texttt{<base>.yaml}, the solid SG YAML documented above, and
\texttt{<base>.msh}, a gmsh v2.2 mesh carrying the material id as both physical
and elementary tag --- and returns the parsed dict. The module defines no
\texttt{\_\_main\_\_} block, so import \texttt{convert} rather than invoking the
module with \texttt{python -m}. Because \texttt{plate\_homo\_2d} performs the
same conversion automatically when handed a \texttt{.sc} path, converting by hand
is only needed when you want to inspect or edit the YAML before homogenizing.

Two failure modes are worth knowing. The tet10 slot convention is remapped
explicitly (four corners, then the six midsides from slots 7--12), and any 3-D
record whose nonzero slot count is not 4, 8 or 10 is refused loudly rather than
converted into phantom nodes.

%========================================================================
\section{The through-thickness \texttt{layup\_db.yaml}}
\label{sec:input-layupdb}

For a plain laminate there is no mesh to write: the 1-D SG is one element stack
through the thickness, and \texttt{layup\_db.yaml} is the whole user input. The
driver \texttt{1d\_sg.py} reads it, generates the SG mesh (\texttt{1dsg.yaml}
plus \texttt{1dsg.png}) through
\texttt{opensg\_solid.rm\_plate\_1D.segment\_plate.plate\_sg\_yaml}, homogenizes
with \texttt{msg\_rm\_plate.rm\_plate\_msg}, and writes
\texttt{layup\_db\_plate\_homo.out}. The generated mesh is shown in
Figure~\ref{fig:input-plate1dsg} and the file's keys in
Table~\ref{tab:input-layupdb}.

\begin{figure}[htbp]
\centering
\includegraphics[width=0.62\textwidth]{\FIGDIR/plate_1dsg.png}
\caption{The generated through-thickness 1-D plate SG, elements coloured by
material: eight thin graphite/epoxy plies bracketing the thick PVC foam core.
The abscissa is $y_3$, running from $-h/2$ to $+h/2$ because
\texttt{fraction: 0.5} put the origin at mid-thickness.}
\label{fig:input-plate1dsg}
\end{figure}

\begin{table}[htbp]
\centering
\small
\caption{Keys of \texttt{layup\_db.yaml} read by \texttt{1d\_sg.py}.}
\label{tab:input-layupdb}
\begin{tabular}{@{}lp{9.8cm}@{}}
\toprule
key & meaning \\
\midrule
\texttt{model} & \texttt{0} prints the classical $6\times6$ ABD, \texttt{1} the
  shear-refined $8\times8$ ABDG. This also names the matrix in the
  \texttt{.out}: \texttt{Classical Plate} or \texttt{Reissner-Mindlin Plate}. \\
\texttt{fraction} & the reference surface as a fraction of the total thickness:
  \texttt{0} = bottom/OML face, \texttt{0.5} = mid-surface, \texttt{1} = top.
  This is the solid-side analogue of the shell YAML's \texttt{reference}. \\
\texttt{mesh.n\_per\_layer} & SG elements per physical ply. \\
\texttt{mesh.elem\_order} & element polynomial order $p$; nodes per element are
  $p+1$, and the engine supports 1 to 4. \texttt{4} --- the five-noded quartic
  --- is the default in the example and represents the whole warping ladder
  exactly ($V_0$ quadratic, $V_1$ cubic, $V_2$ quartic). \\
\texttt{materials} & a mapping from short name to \texttt{E} (3), \texttt{G}
  (3), \texttt{nu} (3) and \texttt{rho}, with an optional \texttt{full\_name}.
  Every number must parse as a float; the example's comment about
  \texttt{128.0e9} is a real PyYAML trap, which is why \texttt{1d\_sg.py}
  coerces with \texttt{float(v)} on read. \\
\texttt{layup} & the stacking sequence, \textbf{bottom ply first}, each entry
  \texttt{\{material, thickness, angle\}}. An optional \texttt{divisions} key is
  read only by the Abaqus 3-D deck generator and ignored by the homogenizer. \\
\bottomrule
\end{tabular}
\end{table}

The Nayak sandwich in the example is a nine-ply $(0/90/0/90/\text{core})_s$ stack
of graphite/epoxy faces on a HEREX C70.130 PVC foam core:

\begin{lstlisting}[style=shellst,caption={The mesh, reference and stacking blocks of layup\_db.yaml.}]
mesh:
  n_per_layer: 1
  elem_order: 4
fraction: 0.5
layup:
  - {material: ge,    thickness: 0.0009525, angle:  0.0}
  - {material: ge,    thickness: 0.0009525, angle: 90.0}
  - {material: ge,    thickness: 0.0009525, angle:  0.0}
  - {material: ge,    thickness: 0.0009525, angle: 90.0}
  - {material: herex, thickness: 0.1447800, angle:  0.0, divisions: 8}
  - {material: ge,    thickness: 0.0009525, angle: 90.0}
  - {material: ge,    thickness: 0.0009525, angle:  0.0}
  - {material: ge,    thickness: 0.0009525, angle: 90.0}
  - {material: ge,    thickness: 0.0009525, angle:  0.0}
\end{lstlisting}

\begin{lstlisting}[style=shellst]
cd examples/OpenSG-solid/1_get_plate_props_from_1DSG
python 1d_sg.py
\end{lstlisting}

\texttt{1d\_sg.py} accepts an alternative database as its one positional
argument, and the output file is named after that database rather than after the
mesh. The \texttt{.out} banner echoes back every number that mattered, which is
the fastest way to confirm the file was read as intended:
\texttt{ OpenSG msg-solid plate model from 1dsg.yaml: 9 plies, h = 0.152400 m,
fraction = 0.5, rho*h = 30.251400 kg/m\^{}2, rows [e11 e22 g12 k11 k22 k12 2g13
2g23]}. With \texttt{model: 1} the matrix printed is the $8\times8$, whose
transverse-shear block for this sandwich is $G_{11} = 7.7010\times10^{6}$ and
$G_{22} = 7.5424\times10^{6}$~N/m, and whose $B$ block is numerically zero
(entries of order $10^{-6}$ against a $10^{8}$ membrane term) as a symmetric
stack requires.

The generated \texttt{1dsg.yaml} records the same decisions in its own header,

\begin{lstlisting}[style=shellst,caption={Header, first node and first elements of the generated 1dsg.yaml (37 node rows and 9 element rows in full).}]
sg: {type: plate_1d, elem_order: 4, n_per_layer: 1, reference_fraction: 0.5,
  thickness: 0.15239999999999998, n_ply: 9}
nodes:
- [-0.07619999999999999]
elements:
- [1, 2, 3, 4, 5]
- [5, 6, 7, 8, 9]
\end{lstlisting}

and \texttt{read\_plate\_sg\_yaml} cross-checks the stored ply thicknesses and
\texttt{reference\_fraction} against the mesh itself, rejecting a hand-edited
file whose header and coordinates disagree instead of silently trusting it.

Note the element rows carefully, because this is the one place where the two
solid routes genuinely differ. \texttt{segment\_plate.plate\_sg\_yaml} numbers
each element's nodes \emph{sequentially}, \texttt{np.arange(p*e, p*e + p + 1)},
which is why the rows read \texttt{[1, 2, 3, 4, 5]} and \texttt{[5, 6, 7, 8, 9]}
--- consecutive integers, adjacent elements sharing their end node.
\texttt{rm\_plate\_msg} then homogenizes that mesh with its own equispaced
Lagrange basis (\texttt{nodes\_xi = np.linspace(-1.0, 1.0, p + 1)},
$\max(3, p+1)$-point Gauss-Legendre), so the \texttt{layup\_db.yaml} route never
touches the general SG engine's element table at all. The ends-first
gmsh/basix line ordering --- \texttt{[seq[0], seq[-1]] + seq[1:-1]} --- belongs
instead to \texttt{opensg\_solid.sg\_mesh.laminate\_to\_sg}, the parity route
that feeds the general engine, and it is that engine's
\texttt{plate\_homo\_2d} which consults \texttt{\_cell\_basis}
(Section~\ref{sec:input-elements}).

%========================================================================
\section{Finite elements used by each route}
\label{sec:input-elements}

The two packages discretize differently, and knowing which element you are
getting explains both the DOF count and the failure modes.

\subsection{Shell route}

The shell SG never assembles a line element directly. The contour is extruded one
quad deep along the beam axis by \texttt{sg\_homo.\_strip}, with extrusion length
$h$ equal to the mean contour element length, giving quads
\texttt{[a, b, m+b, m+a]}; the top node row is then DOF-mapped onto the bottom
row (\texttt{dof\_map = np.concatenate([np.arange(m), np.arange(m)])}) so the
field is exactly span-invariant. That is the operator's own prismatic reduction.
The element is therefore a 4-node bilinear quad integrated at $2\times2$ Gauss
points ($\pm 1/\sqrt{3}$ in each direction), and the 1-D cross-section problem is
recovered as its span-invariant limit. Two element families are built on that
quad, and they differ only in what they do with the drilling rotation
(Table~\ref{tab:input-shell-elements}).

\begin{table}[htbp]
\centering
\small
\caption{The two shell element families and how each handles the drilling
rotation.}
\label{tab:input-shell-elements}
\begin{tabular}{@{}lccp{5.6cm}@{}}
\toprule
entry point & DOF/node & DOF/elem & drilling \\
\midrule
\texttt{sg\_homo.ring\_indep} (production)
  & 6: \texttt{[w1, w2, w3, om1, om2, om3]} & 24
  & independent nodal $\omega_3$, tied back by a finite drilling residual
    (\texttt{NDOF6 = 6}) \\
\texttt{sg\_assembly.ring\_general}
  & 5: \texttt{[w1, w2, w3, om1, om2]} & 20
  & eliminated algebraically through
    $\omega_3 = S/(2C_{33}) - (C_{3b}/C_{33})\,\omega_b$ (\texttt{NDOF = 5}) \\
\bottomrule
\end{tabular}
\end{table}

The independent-drilling element exists because the elimination divides by
$C_{33} = n \cdot b_3$, which vanishes on a flat wall. In the 6-DOF element
$\omega_3$ is a genuine nodal DOF and the in-plane symmetry that defined it is
re-imposed in its finite, undivided form as the drilling residual
\[
DR = C_{33}\,\omega_3 + C_{3b}\,\omega_b - S/2 ,
\]
which equals $C_{33}(\omega_3 - \omega_3^{\text{eliminated}})$ and stays finite
even at $C_{33} = 0$. It is enforced by an element-wise Lagrange multiplier
(\texttt{lam\_space="elem"}, the production setting) or by a penalty. Transverse
shear uses MITC assumed-strain tying in the Dvorkin--Bathe sense: the production
ring scheme is \texttt{shear="mitc4\_g23"}, which ties only $\gamma_{23}$,
because under span invariance $\gamma_{13}$ is algebraic in the directors and
carries no fluctuation gradient.

The 3-D shell SG uses the same quad element, with triangles entering as quads
whose last node is repeated.

\subsection{Solid route}

\texttt{opensg\_solid.sg\_mesh.\_cell\_basis} maps the SG dimension and the node
count per element onto a basix cell type and Lagrange degree. This table
\emph{is} the list of supported elements, and it is consulted by exactly one
function: \texttt{sg\_homo.plate\_homo\_2d}, the general SG engine entry point
(and therefore also by anything that reaches it, including a \texttt{.sc} handed
to it directly and the \texttt{laminate\_to\_sg} parity route). The supported
combinations are given in Table~\ref{tab:input-cellbasis}.

\begin{table}[htbp]
\centering
\small
\caption{The \texttt{\_cell\_basis} table: every element the general solid SG
engine supports.}
\label{tab:input-cellbasis}
\begin{tabular}{@{}cccc@{}}
\toprule
\texttt{dim} & nodes/elem & basix cell & degree \\
\midrule
1 & 2 & \texttt{interval} & 1 \\
1 & 3 & \texttt{interval} & 2 \\
1 & 4 & \texttt{interval} & 3 \\
1 & 5 & \texttt{interval} & 4 \\
2 & 3 & \texttt{triangle} & 1 \\
2 & 4 & \texttt{quadrilateral} & 1 \\
3 & 4 & \texttt{tetrahedron} & 1 \\
3 & 8 & \texttt{hexahedron} & 1 \\
3 & 10 & \texttt{tetrahedron} & 2 \\
\bottomrule
\end{tabular}
\end{table}

Anything else raises \texttt{unsupported <dim>D element with <n> nodes}. Basis
functions are equispaced Lagrange (\texttt{LagrangeVariant.equispaced}) with
quadrature degree 6 for degree $>1$ and 2 otherwise. Two connectivity
conventions are applied on the way in: \texttt{\_to\_basix\_order} permutes quad4
to \texttt{[0, 1, 3, 2]} and hex8 to \texttt{[0, 1, 3, 2, 4, 5, 7, 6]}, while
triangles, tetrahedra and intervals pass through unchanged.

The through-thickness laminate of Section~\ref{sec:input-layupdb} does
\emph{not} go through this table. Its route ---
\texttt{plate\_sg\_yaml} to \texttt{rm\_plate\_msg} --- carries its own
equispaced Lagrange basis and its own sequential node numbering, as described
there. Only \texttt{laminate\_to\_sg}, the alternative route that packages a
laminate as a generic 1-D SG dict for \texttt{plate\_homo\_2d}, writes the
ends-first gmsh/basix line ordering that the \texttt{\_cell\_basis} elements
expect.

%========================================================================
\section{Sign and order conventions}
\label{sec:input-conventions}

Three orderings appear across the codebase. Every one of them is stated in a
module docstring rather than inferred, and every one of them is a common source
of a plausible-looking but wrong answer.

\subsection{Solid (3-D) Voigt order}

The solid order is the SwiftComp order
\[
[\,11\;\;22\;\;33\;\;23\;\;13\;\;12\,]
\]
with \textbf{engineering shears}. The compliance therefore carries
$S_{44} = 1/G_{23}$, $S_{55} = 1/G_{13}$, $S_{66} = 1/G_{12}$, and the shear rows
of a recovered strain are $2\gamma$, not $\gamma$. This is the order of
\texttt{write\_sc\_K}'s output, of \texttt{build\_stiffness\_6x6}, and of every
\texttt{type: 2} material \texttt{C} block in a solid SG YAML.

\subsection{Beam order}

The beam order is
\[
[\,\varepsilon_{11}\;\;\gamma_{12}\;\;\gamma_{13}\;\;\kappa_1\;\;\kappa_2\;\;
\kappa_3\,],
\]
whose diagonal is the familiar VABS ordering
\[
[\,EA\;\;GA_2\;\;GA_3\;\;GJ\;\;EI_2\;\;EI_3\,],
\]
the \texttt{LBL} list used by the shell drivers. \textbf{Axis 1 is the beam
axis}; axes 2 and 3 are the cross-section axes.

The solver indexes global axes as $(1,2,3) = (\text{axial},\,y_2,\,y_3)$ while an
SG YAML writes each 3-vector as $(y_2,\,y_3,\,\text{axial})$. That is why
\texttt{opensg\_solid.sg\_materials.elem\_rotation\_from\_yaml} cycles
\[
[a,\,b,\,c] \;\longrightarrow\; [c,\,a,\,b]
\]
(implemented as \texttt{o[:, :, [2, 0, 1]]}) before a YAML orientation can be
used as a direction-cosine matrix in the solid beam path. Skipping that cycle
points $e_1$ along $y_3$: the run still completes and still looks plausible, but
it destroys the section's symmetry, and a square tube comes out with
$EI_2 \neq EI_3$.

\subsection{Plate order}

The plate order is $[\,N_{11}\;N_{22}\;N_{12}\;M_{11}\;M_{22}\;M_{12}\,]$ for the
classical $6\times6$, conjugate to the strain vector
$[\varepsilon_{11}\;\varepsilon_{22}\;2\varepsilon_{12}\;\kappa_{11}\;
\kappa_{22}\;\kappa_{12}]$. The Reissner--Mindlin $8\times8$ appends the two
transverse shears, giving the rows
\texttt{[e11 e22 g12 k11 k22 k12 2g13 2g23]} that both
\texttt{write\_abdg\_out} and \texttt{1d\_sg.py} print in their banners, and the
block structure
\[
\mathbf{ABDG} =
\begin{bmatrix}
\mathbf{A} & \mathbf{B} & \mathbf{0} \\
\mathbf{B} & \mathbf{D} & \mathbf{0} \\
\mathbf{0} & \mathbf{0} & \mathbf{G}
\end{bmatrix}.
\]

For a plate SG the \textbf{last SG coordinate is the thickness direction} --- the
$y_3$ of the solid Voigt order above. With \texttt{dim: 2} that is node
coordinate index 1, since only the first \texttt{dim} coordinates are read; for a
through-thickness 1-D SG it is the single coordinate. This is why the RHC cell
reports $\omega = 35.248$, the span of coordinate 0 (the in-plane periodic
direction), and not its larger extent in coordinate 1.

\subsection{Changing the reference surface}

Moving the reference surface by $z_0$ maps the membrane strain
$\varepsilon$ to $\varepsilon + z_0\,\kappa$, so with
$\mathbf{T} = \begin{bmatrix} \mathbf{I} & z_0\mathbf{I} \\ \mathbf{0} &
\mathbf{I}\end{bmatrix}$ the plate law transforms as
\[
\mathbf{ABD}_{\text{new}} = \mathbf{T}^{-\mathsf{T}}\,\mathbf{ABD}\,
\mathbf{T}^{-1},
\]
which is exactly what \texttt{fe\_jax.msg\_materials.shift\_abd\_reference(ABD,
z0)} computes. Use it rather than reversing the layup: reversal flips $e_3$ and
breaks agreement with the material orientation. The transverse-shear block
$\mathbf{G}$ is reference-independent, which is why
\texttt{compute\_ABD\_matrix} can shift the $6\times6$ and leave the $8\times8$'s
$\mathbf{G}$ untouched.

%========================================================================
\section{Material orientation}
\label{sec:input-orientation}

Each row of \texttt{elementOrientations} is nine direction cosines forming three
stacked unit vectors,

\begin{lstlisting}[style=shellst]
[ e1x e1y e1z | e2x e2y e2z | e3x e3y e3z ]
\end{lstlisting}

with the roles fixed across both shell dialects:

\begin{itemize}
\item $e_1$ is the \textbf{beam axis} --- the prismatic direction, which for a
  plane cross-section meshed in $(y_2, y_3)$ is out of plane, so
  \texttt{e1 = [0, 0, 1]}.
\item $e_2$ is the \textbf{in-plane ply-flow direction}, the wall tangent along
  the contour.
\item $e_3 = e_1 \times e_2$ is the \textbf{wall normal}, in plane for a
  cross-section SG, and the direction the layup stacks along from the reference
  surface.
\end{itemize}

The triad must be orthonormal and right-handed.
\texttt{opensg\_shell.sg\_mesh.frame\_report(nodes, cells, e1, e2, e3)} checks
this numerically. It \textbf{returns} the pair \texttt{(ok, txt)} and prints
nothing --- callers such as \texttt{sg\_mesh.extract} do the printing --- and its
one-line summary ends in \texttt{OK} or \texttt{CHECK}. Its pass/fail gate is

\[
\bigl|\,\overline{|e_i|} - 1\,\bigr| < 10^{-6}\ (i = 1,2,3), \qquad
\max_{i\neq j}\bigl|e_i \cdot e_j\bigr| < 10^{-6},
\]
\[
\min\bigl(e_1 \times e_2 \cdot e_3\bigr) > 0.99, \qquad
\min\bigl(e_1 \cdot \hat{x}\bigr) > 0.999 .
\]

The first three are universal; the fourth --- $e_1$ axial --- is
cylinder-specific, as the docstring says. The report also prints a fifth number,
the mean of $e_3 \cdot (-\hat{r})$, i.e.\ how nearly $e_3$ is the inward radial
normal, but that value is \textbf{reported and not gated}: it appears in the text
so you can read it, and a poor value alone will not turn an \texttt{OK} into a
\texttt{CHECK}.

Because any mesh handling can renumber nodes and cells, an orientation PNG is a
compulsory deliverable of every OpenSG run. It is emitted by the \emph{example
driver} calling \texttt{opensg\_shell.auto\_emit} --- see the last line of
\texttt{beam\_homo\_shell.py} --- and not by any homogenization entry point:
\texttt{build\_rm\_bundle}, \texttt{shell\_sg3d} and \texttt{plate\_homo\_2d}
never call it. \texttt{auto\_emit} is cached once per \texttt{(shell, solid)}
path per process and swallows any exception with a
\texttt{[orient\_plot] skipped} message, so it never raises and is safe to call
from inside a compute path; a new driver of your own should call it explicitly.

The colour convention of \texttt{fe\_jax/orient\_plot.py} is \textbf{$e_2$ blue,
$e_3$ green}, drawn at element centroids over a light mesh; $e_1$ is out of the
page and so not drawn. (The separate 3-D plotters in \texttt{sg\_mesh.py} use
red/blue/black for $e_1$/$e_2$/$e_3$, since there $e_1$ is visible.) The
canonical example is the IEA-22 $r/R = 0.2$ station of
Figure~\ref{fig:input-orient}.

\begin{figure}[htbp]
\centering
\includegraphics[width=\textwidth]{\FIGDIR/shell_xsec_orient.png}
\caption{Material orientation of the IEA-22 $r/R = 0.2$ shell cross-section:
$e_2$ in blue (in-plane ply direction), $e_3$ in green (wall normal). This is
\texttt{iea\_s10\_orient.png}, written by the \texttt{auto\_emit} call at the end
of \texttt{beam\_homo\_shell.py}.}
\label{fig:input-orient}
\end{figure}

Read it as a physical check, not decoration. On this station the green $e_3$
arrows point inward everywhere on the skin, which is what an outward-first layup
stacked from the mid-surface requires, and the blue $e_2$ arrows are tangent to
every wall --- running consistently around the closed skin loop and along each of
the three shear webs. An $e_3$ that flips between neighbouring skin elements, or
a web whose $e_2$ opposes its neighbours, means the layup is being stacked in two
directions at once: visible here in seconds, and invisible in the resulting
$6\times6$.

\subsection{Where the ply angle lives}

The ply angle is \textbf{not} baked into $e_1$. It lives in the third column of
each \texttt{layup} row (shell dialect) or in the \texttt{angles} argument (solid
dialect), and it rotates the material about $e_3$, away from $e_1$. The rotation
\texttt{rotation\_6x6(theta)} leaves the Voigt-33 row untouched and mixes the
11/22/12 rows, and the reduced in-plane stiffness built by
\texttt{msg\_transverse\_shear.\_ply\_Q\_and\_G} is the textbook
\[
\bar{Q}_{11} = Q_{11}\cos^4\theta + 2(Q_{12} + 2Q_{66})\sin^2\theta\cos^2\theta
  + Q_{22}\sin^4\theta ,
\]
so \texttt{angle: 0.0} places the fibre along $e_1$ --- the beam axis. A
wind-blade unidirectional spar cap therefore carries \texttt{angle: 0.0}, and a
$\pm45$ ply carries \texttt{45.0} and \texttt{-45.0}.

The governing repository rule is
\texttt{Rules/orientation\_e1\_out\_of\_plane.md}, which also records the failure
mode: a YAML whose $e_1$ already contains the fibre rotation (so
$|e_{1z}| < 1$) is a \emph{fibre} frame, and combining it with a non-zero
\texttt{angle} applies the rotation twice. The rule supplies a two-line check
worth running on any YAML you did not write yourself:

\begin{lstlisting}[style=opensg,caption={The orientation precheck from Rules/orientation\_e1\_out\_of\_plane.md.}]
ori = np.array(d["elementOrientations"], float)
e1, e3 = ori[:, 0:3], ori[:, 6:9]
assert np.allclose(np.abs(e1[:, 2]), 1.0)   # e1 out of plane
assert np.allclose(e3[:, 2], 0.0)           # e3 in plane
\end{lstlisting}

If the first assertion fails, the file carries a fibre frame; either switch to
the frame variant built with the axial rotation only, or set the fibre angle to
zero downstream so the rotation is applied exactly once.
